\documentclass[12pt]{article}
\usepackage[utf8]{inputenc}
\usepackage[T1]{fontenc}
\usepackage{hyperref}
\title{Bactrian Language Resources for K-12 Curriculum Development}
\author{Compiled by Claude Sonnet 5.5 (claude-sonnet-5-5) for Andr\'as Kornai}
\date{}
\begin{document}
\maketitle

\section{Introduction}
Bactrian (xbc) is an extinct Middle Iranian language of ancient Bactria (today northern Afghanistan). It is the only Middle Iranian language written in a Greek-based alphabet. The best-known text is the Rabatak inscription, found in 1993 and dating from the Kushan period (2nd century CE). Before 2000 knowledge rested mostly on a few inscriptions plus coin and seal legends. Documents from northern Afghanistan were then published by Nicholas Sims-Williams, who has been the main scholar editing them. There is no native-speaker community, and I found no school materials.

\section{Primary Educational Resources}
None found. No primers, textbooks or school programmes for Bactrian appeared in the search results.

\section{Linguistic and Reference Materials}
Sims-Williams, \emph{Bactrian Documents from Northern Afghanistan} (Studies in the Khalili Collection, vol.\ III): Part I (legal and economic documents, revised edition 2012), Part II (letters and Buddhist texts, 2007, with an outline grammar and glossary, listed as out of print) and Part III (plates). The Khalili Collections site lists these, with a free PDF of Part I after registration according to the search summary: \url{https://www.khalilicollections.org/?p=3958}. Overview articles: \url{https://en.wikipedia.org/wiki/Rabatak_inscription}, \url{https://www.Wikipedia.com/wiki/Bactrian_language}, and the Omniglot page on the script: \url{https://omniglot.com/writing/bactrian.htm}. An article on the Rabatak inscription is at \url{https://tesl.shirazu.ac.ir/article_255.html}. I found no dictionary in the results.

\section{Religious and Community Resources}
None found. The documents include Buddhist texts (per the Khalili listing), but I saw no URL for a specific edition beyond the volumes above.

\section{Digital and Audio Resources}
No online searchable corpus was found. The Khalili Collections pages are the closest digital access point. Some documents are held in the British Academy collections: \url{https://collections.thebritishacademy.ac.uk/object-bab-12729}.

\section{Curriculum Development Context}
A K-12 use would be a heritage or classical-languages enrichment unit: the Greek-derived script, the Kushan history, and the primary documents in translation. Teachers would depend on scholarly editions, since no pedagogical materials were found.
\end{document}